% Options for packages loaded elsewhere
\PassOptionsToPackage{unicode}{hyperref}
\PassOptionsToPackage{hyphens}{url}
%
\documentclass[
  a4paper,
  oneside,
  open=any]{scrbook}

\usepackage{amsmath,amssymb}
\usepackage{iftex}
\ifPDFTeX
  \usepackage[T1]{fontenc}
  \usepackage[utf8]{inputenc}
  \usepackage{textcomp} % provide euro and other symbols
\else % if luatex or xetex
  \usepackage{unicode-math}
  \defaultfontfeatures{Scale=MatchLowercase}
  \defaultfontfeatures[\rmfamily]{Ligatures=TeX,Scale=1}
\fi
\usepackage{lmodern}
\ifPDFTeX\else  
    % xetex/luatex font selection
    \setmainfont[]{Public Sans}
\fi
% Use upquote if available, for straight quotes in verbatim environments
\IfFileExists{upquote.sty}{\usepackage{upquote}}{}
\IfFileExists{microtype.sty}{% use microtype if available
  \usepackage[]{microtype}
  \UseMicrotypeSet[protrusion]{basicmath} % disable protrusion for tt fonts
}{}
\makeatletter
\@ifundefined{KOMAClassName}{% if non-KOMA class
  \IfFileExists{parskip.sty}{%
    \usepackage{parskip}
  }{% else
    \setlength{\parindent}{0pt}
    \setlength{\parskip}{6pt plus 2pt minus 1pt}}
}{% if KOMA class
  \KOMAoptions{parskip=half}}
\makeatother
\usepackage{xcolor}
\usepackage[top=30mm,left=25mm,right=25mm,bottom=30mm]{geometry}
\setlength{\emergencystretch}{3em} % prevent overfull lines
\setcounter{secnumdepth}{-\maxdimen} % remove section numbering
% Make \paragraph and \subparagraph free-standing
\makeatletter
\ifx\paragraph\undefined\else
  \let\oldparagraph\paragraph
  \renewcommand{\paragraph}{
    \@ifstar
      \xxxParagraphStar
      \xxxParagraphNoStar
  }
  \newcommand{\xxxParagraphStar}[1]{\oldparagraph*{#1}\mbox{}}
  \newcommand{\xxxParagraphNoStar}[1]{\oldparagraph{#1}\mbox{}}
\fi
\ifx\subparagraph\undefined\else
  \let\oldsubparagraph\subparagraph
  \renewcommand{\subparagraph}{
    \@ifstar
      \xxxSubParagraphStar
      \xxxSubParagraphNoStar
  }
  \newcommand{\xxxSubParagraphStar}[1]{\oldsubparagraph*{#1}\mbox{}}
  \newcommand{\xxxSubParagraphNoStar}[1]{\oldsubparagraph{#1}\mbox{}}
\fi
\makeatother

\usepackage{color}
\usepackage{fancyvrb}
\newcommand{\VerbBar}{|}
\newcommand{\VERB}{\Verb[commandchars=\\\{\}]}
\DefineVerbatimEnvironment{Highlighting}{Verbatim}{commandchars=\\\{\}}
% Add ',fontsize=\small' for more characters per line
\usepackage{framed}
\definecolor{shadecolor}{RGB}{241,243,245}
\newenvironment{Shaded}{\begin{snugshade}}{\end{snugshade}}
\newcommand{\AlertTok}[1]{\textcolor[rgb]{0.68,0.00,0.00}{#1}}
\newcommand{\AnnotationTok}[1]{\textcolor[rgb]{0.37,0.37,0.37}{#1}}
\newcommand{\AttributeTok}[1]{\textcolor[rgb]{0.40,0.45,0.13}{#1}}
\newcommand{\BaseNTok}[1]{\textcolor[rgb]{0.68,0.00,0.00}{#1}}
\newcommand{\BuiltInTok}[1]{\textcolor[rgb]{0.00,0.23,0.31}{#1}}
\newcommand{\CharTok}[1]{\textcolor[rgb]{0.13,0.47,0.30}{#1}}
\newcommand{\CommentTok}[1]{\textcolor[rgb]{0.37,0.37,0.37}{#1}}
\newcommand{\CommentVarTok}[1]{\textcolor[rgb]{0.37,0.37,0.37}{\textit{#1}}}
\newcommand{\ConstantTok}[1]{\textcolor[rgb]{0.56,0.35,0.01}{#1}}
\newcommand{\ControlFlowTok}[1]{\textcolor[rgb]{0.00,0.23,0.31}{\textbf{#1}}}
\newcommand{\DataTypeTok}[1]{\textcolor[rgb]{0.68,0.00,0.00}{#1}}
\newcommand{\DecValTok}[1]{\textcolor[rgb]{0.68,0.00,0.00}{#1}}
\newcommand{\DocumentationTok}[1]{\textcolor[rgb]{0.37,0.37,0.37}{\textit{#1}}}
\newcommand{\ErrorTok}[1]{\textcolor[rgb]{0.68,0.00,0.00}{#1}}
\newcommand{\ExtensionTok}[1]{\textcolor[rgb]{0.00,0.23,0.31}{#1}}
\newcommand{\FloatTok}[1]{\textcolor[rgb]{0.68,0.00,0.00}{#1}}
\newcommand{\FunctionTok}[1]{\textcolor[rgb]{0.28,0.35,0.67}{#1}}
\newcommand{\ImportTok}[1]{\textcolor[rgb]{0.00,0.46,0.62}{#1}}
\newcommand{\InformationTok}[1]{\textcolor[rgb]{0.37,0.37,0.37}{#1}}
\newcommand{\KeywordTok}[1]{\textcolor[rgb]{0.00,0.23,0.31}{\textbf{#1}}}
\newcommand{\NormalTok}[1]{\textcolor[rgb]{0.00,0.23,0.31}{#1}}
\newcommand{\OperatorTok}[1]{\textcolor[rgb]{0.37,0.37,0.37}{#1}}
\newcommand{\OtherTok}[1]{\textcolor[rgb]{0.00,0.23,0.31}{#1}}
\newcommand{\PreprocessorTok}[1]{\textcolor[rgb]{0.68,0.00,0.00}{#1}}
\newcommand{\RegionMarkerTok}[1]{\textcolor[rgb]{0.00,0.23,0.31}{#1}}
\newcommand{\SpecialCharTok}[1]{\textcolor[rgb]{0.37,0.37,0.37}{#1}}
\newcommand{\SpecialStringTok}[1]{\textcolor[rgb]{0.13,0.47,0.30}{#1}}
\newcommand{\StringTok}[1]{\textcolor[rgb]{0.13,0.47,0.30}{#1}}
\newcommand{\VariableTok}[1]{\textcolor[rgb]{0.07,0.07,0.07}{#1}}
\newcommand{\VerbatimStringTok}[1]{\textcolor[rgb]{0.13,0.47,0.30}{#1}}
\newcommand{\WarningTok}[1]{\textcolor[rgb]{0.37,0.37,0.37}{\textit{#1}}}

\providecommand{\tightlist}{%
  \setlength{\itemsep}{0pt}\setlength{\parskip}{0pt}}\usepackage{longtable,booktabs,array}
\usepackage{calc} % for calculating minipage widths
% Correct order of tables after \paragraph or \subparagraph
\usepackage{etoolbox}
\makeatletter
\patchcmd\longtable{\par}{\if@noskipsec\mbox{}\fi\par}{}{}
\makeatother
% Allow footnotes in longtable head/foot
\IfFileExists{footnotehyper.sty}{\usepackage{footnotehyper}}{\usepackage{footnote}}
\makesavenoteenv{longtable}
\usepackage{graphicx}
\makeatletter
\def\maxwidth{\ifdim\Gin@nat@width>\linewidth\linewidth\else\Gin@nat@width\fi}
\def\maxheight{\ifdim\Gin@nat@height>\textheight\textheight\else\Gin@nat@height\fi}
\makeatother
% Scale images if necessary, so that they will not overflow the page
% margins by default, and it is still possible to overwrite the defaults
% using explicit options in \includegraphics[width, height, ...]{}
\setkeys{Gin}{width=\maxwidth,height=\maxheight,keepaspectratio}
% Set default figure placement to htbp
\makeatletter
\def\fps@figure{htbp}
\makeatother
% definitions for citeproc citations
\NewDocumentCommand\citeproctext{}{}
\NewDocumentCommand\citeproc{mm}{%
  \begingroup\def\citeproctext{#2}\cite{#1}\endgroup}
\makeatletter
 % allow citations to break across lines
 \let\@cite@ofmt\@firstofone
 % avoid brackets around text for \cite:
 \def\@biblabel#1{}
 \def\@cite#1#2{{#1\if@tempswa , #2\fi}}
\makeatother
\newlength{\cslhangindent}
\setlength{\cslhangindent}{1.5em}
\newlength{\csllabelwidth}
\setlength{\csllabelwidth}{3em}
\newenvironment{CSLReferences}[2] % #1 hanging-indent, #2 entry-spacing
 {\begin{list}{}{%
  \setlength{\itemindent}{0pt}
  \setlength{\leftmargin}{0pt}
  \setlength{\parsep}{0pt}
  % turn on hanging indent if param 1 is 1
  \ifodd #1
   \setlength{\leftmargin}{\cslhangindent}
   \setlength{\itemindent}{-1\cslhangindent}
  \fi
  % set entry spacing
  \setlength{\itemsep}{#2\baselineskip}}}
 {\end{list}}
\usepackage{calc}
\newcommand{\CSLBlock}[1]{\hfill\break\parbox[t]{\linewidth}{\strut\ignorespaces#1\strut}}
\newcommand{\CSLLeftMargin}[1]{\parbox[t]{\csllabelwidth}{\strut#1\strut}}
\newcommand{\CSLRightInline}[1]{\parbox[t]{\linewidth - \csllabelwidth}{\strut#1\strut}}
\newcommand{\CSLIndent}[1]{\hspace{\cslhangindent}#1}

\usepackage{fontspec}
\setmainfont{KoPubWorldDotum Light}  % Set main body font with fallback
\setsansfont{KoPubWorldBatang Bold}       % Sans-serif font with fallback
\setmonofont{D2Coding}  % Monospaced font with fallback
\usepackage{graphicx}
\usepackage{eso-pic}
\usepackage{lastpage}
\usepackage{wallpaper} % for the background image on title page
\usepackage{fontawesome5}
\usepackage{xcolor}
\definecolor{orcidlogocol}{rgb}{0.65, 0.807, 0.223}
\definecolor{anugoldtint}{HTML}{F5EDDE}
\definecolor{anugold}{HTML}{BE830E}
\newcommand{\orcid}[1]{\href{https://orcid.org/#1}{\textcolor{orcidlogocol}{\faOrcid} #1}}

\usepackage[manualmark]{scrlayer-scrpage}
\clearpairofpagestyles
\ihead{}
\ofoot{Page \thepage\ of \pageref*{LastPage}}
\ifoot{The Korea University\\{\footnotesize TEQSA Provider ID: PRV12002 (Korea University) | CRICOS Provider Code: 00120C}}
\addtokomafont{pageheadfoot}{\upshape}

\usepackage{multicol}

\usepackage[most]{tcolorbox}
\newtcolorbox{titlepagebox}{colback=anugoldtint,colframe=white}
\newtcbox{\inlinebox}[1][]{enhanced,
 box align=base,
 nobeforeafter,
 colback=white,
 colframe=gray,
 size=small,
 left=0pt,
 right=0pt,
 boxsep=2pt,
 #1}

% \usepackage{fancyhdr}
% % for title page
% \fancypagestyle{plain}{%
%   \renewcommand{\headrulewidth}{0pt}%
%   \fancyhf{}%
%   \fancyhf[rf]{Page \thepage\ of \pageref*{LastPage}}
% \fancyhf[lf]{The Australian National University\\{\footnotesize TEQSA Provider ID: PRV12002 (Australian University) | CRICOS Provider Code: 00120C}}
% }
% % for other pages
% \pagestyle{fancy}
% \fancyhf{}
% \fancyhf[rf]{Page \thepage\ of \pageref*{LastPage}}
% \fancyhf[lf]{The Australian National University\\{\footnotesize TEQSA Provider ID: PRV12002 (Australian University) | CRICOS Provider Code: 00120C}}
% \renewcommand{\headrulewidth}{0pt} % removes horizontal header line

%\usepackage{titling}
%\pretitle{% add some rules
%   \Huge\bfseries
%}%, make the fonts bigger, make the title (only) bold
%\posttitle{%
% \vskip .75em plus .25em minus .25em% increase the vertical spacing a bit, make this particular glue stretchier
%}

% ensure chapter starts on the same page
\makeatletter
\patchcmd{\scr@startchapter}{\if@openright\cleardoublepage\else\clearpage\fi}{}{}{}
\makeatother
\makeatletter
\@ifpackageloaded{caption}{}{\usepackage{caption}}
\AtBeginDocument{%
\ifdefined\contentsname
  \renewcommand*\contentsname{Table of contents}
\else
  \newcommand\contentsname{Table of contents}
\fi
\ifdefined\listfigurename
  \renewcommand*\listfigurename{List of Figures}
\else
  \newcommand\listfigurename{List of Figures}
\fi
\ifdefined\listtablename
  \renewcommand*\listtablename{List of Tables}
\else
  \newcommand\listtablename{List of Tables}
\fi
\ifdefined\figurename
  \renewcommand*\figurename{Figure}
\else
  \newcommand\figurename{Figure}
\fi
\ifdefined\tablename
  \renewcommand*\tablename{Table}
\else
  \newcommand\tablename{Table}
\fi
}
\@ifpackageloaded{float}{}{\usepackage{float}}
\floatstyle{ruled}
\@ifundefined{c@chapter}{\newfloat{codelisting}{h}{lop}}{\newfloat{codelisting}{h}{lop}[chapter]}
\floatname{codelisting}{Listing}
\newcommand*\listoflistings{\listof{codelisting}{List of Listings}}
\makeatother
\makeatletter
\makeatother
\makeatletter
\@ifpackageloaded{caption}{}{\usepackage{caption}}
\@ifpackageloaded{subcaption}{}{\usepackage{subcaption}}
\makeatother

\ifLuaTeX
  \usepackage{selnolig}  % disable illegal ligatures
\fi
\usepackage{bookmark}

\IfFileExists{xurl.sty}{\usepackage{xurl}}{} % add URL line breaks if available
\urlstyle{same} % disable monospaced font for URLs
\hypersetup{
  pdftitle={국세청 천분위 소득자료(2012-2023)},
  pdfauthor={박영삼; Prof.~Florence Nightingale; Ronald Fisher; Chat GPT},
  pdfkeywords={국세청 천분위 소득자료, 개인소득분배, 불평등},
  hidelinks,
  pdfcreator={LaTeX via pandoc}}


\title{국세청 천분위 소득자료(2012-2023)}
\usepackage{etoolbox}
\makeatletter
\providecommand{\subtitle}[1]{% add subtitle to \maketitle
  \apptocmd{\@title}{\par {\large #1 \par}}{}{}
}
\makeatother
\subtitle{- 통합소득/근로소득/종합소득 천분위 자료 분석 결과 -}
\author{박영삼 \and Prof.~Florence Nightingale \and Ronald
Fisher \and Chat GPT}
\date{25th March 2025}

\begin{document}
% Inspired by https://github.com/nmfs-opensci/quarto_titlepages_v1
  \begin{frontmatter}
  \begin{titlepage}
  % This is a combination of Pandoc templating and LaTeX
  % Pandoc templating https://pandoc.org/MANUAL.html#templates
  % See the README for help

  \newgeometry{top=3in,bottom=1in,right=1in,left=1in}
  \begin{minipage}[b][\textheight][s]{\textwidth}
  \raggedright

  % background image
  \ThisULCornerWallPaper{1}{\_extensions/anu-report/assets/images/coverpage.jpg}

  % Title and subtitle
  {\Huge\bfseries{국세청 천분위 소득자료(2012-2023)}}\\[1\baselineskip]
  {\LARGE{- 통합소득/근로소득/종합소득 천분위 자료 분석 결과
  -}}\\[1\baselineskip]
      \inlinebox{\scriptsize\textcolor{gray}{\MakeUppercase{소득분배}}}
      \inlinebox{\scriptsize\textcolor{gray}{\MakeUppercase{개인소득}}}
      \inlinebox{\scriptsize\textcolor{gray}{\MakeUppercase{국세통계}}}
    \vspace{3mm}

      This report template is available at
      \href{https://github.com/anuopensci/quarto-anu-report}{github.com/anuopensci/quarto-anu-report}.

  \begin{titlepagebox}



  % Authors
  \textcolor{anugold}{\MakeUppercase{Authors}}
  \vspace{-3mm}
  \begin{multicols}{2}
      \begin{minipage}{\columnwidth}
      \raggedright
      \normalfont
      {\textbf{박영삼}}\\
          
      {\itshape{고려대 노동문제연구소 노동데이터센터}}\\
          \faIcon[regular]{envelope} youngsampk@gmail.com\\
      
      \orcid{0000-0001-8771-9509}\\
      \vspace{4mm}
      \end{minipage}

      \begin{minipage}{\columnwidth}
      \raggedright
      \normalfont
      {\textbf{Prof.~Florence Nightingale}}\\
          {Research School of Visualiation}\\
      {\itshape{Australian National University}}\\
          \faIcon[regular]{envelope} florence@anu.edu.au\\
      
      
      \vspace{4mm}
      \end{minipage}

      \begin{minipage}{\columnwidth}
      \raggedright
      \normalfont
      {\textbf{Ronald Fisher}}\\
          {Frequentist Statistics Institute}\\
      {\itshape{Australian National University}}\\
          \faIcon[regular]{envelope} ronald@anu.edu.au\\
      
      
      \vspace{4mm}
      \end{minipage}

      \begin{minipage}{\columnwidth}
      \raggedright
      \normalfont
      {\textbf{Chat GPT}}\\
          
      {\itshape{OpenAI}}\\
          
      
      
      \vspace{4mm}
      \end{minipage}

  \end{multicols}

  \vspace{-4mm}

  % Date %

  % FIXME: \bfseries doesn't seem to be working
  {\textcolor{anugold}{\MakeUppercase{Date}}}\\ 25th March 2025
  {\itshape{(Last modified: 23rd March 2025)}}
  \vspace{3mm}

  \textcolor{anugold}{\MakeUppercase{Executive Summary}}\\
  The report delineates key principles, such as version control and data
  sharing, and introduces practical tools like RMarkdown, Jupyter
  Notebooks, and Docker. Reproducible research significantly boosts the
  credibility of scientific findings, fostering collaboration, iterative
  improvements, and efficient communication. The document addresses
  challenges like data privacy and infrastructure constraints, proposing
  solutions such as anonymization and institutional support. Recognizing
  the transformative impact of reproducibility, the report concludes
  that its adoption contributes to building a culture of transparency
  and accountability in scientific endeavors, laying the groundwork for
  a more robust foundation in advancing knowledge. Recommendations
  include providing training resources, support infrastructure, and
  guidelines to facilitate the integration of reproducible research
  practices into standard workflows.
  \vspace{3mm}

  {\textcolor{anugold}{\MakeUppercase{Keywords}}}\\
  국세청 천분위 소득자료, 개인소득분배, 불평등

  \end{titlepagebox}

  \vspace{2\baselineskip}





  %use \vfill instead to get the space to fill flexibly
  %\vspace{0.25\textheight} % Whitespace between the title block and the publisher

  \vfill

  %%%%%% Cover image


  % Whitespace between the title block and the tagline
  \vspace{1\baselineskip}

  The Korea University\\{\footnotesize TEQSA Provider ID: PRV12002 (Korea University) | CRICOS Provider Code: 00120C}

  \end{minipage}
  \restoregeometry

  \end{titlepage}
  \end{frontmatter}

\restoregeometry

\cleardoublepage
\AddToShipoutPicture{\includegraphics[width=\paperwidth,height=\paperheight]{\_extensions/anu-report/assets/images/page-background.jpg}}


\floatplacement{table}{H}
\renewcommand*\contentsname{Table of contents}
{
\setcounter{tocdepth}{2}
\tableofcontents
}

\mainmatter
\section{Introduction}\label{introduction}

Reproducible research is an essential paradigm that promotes the idea
that scientific investigations should be transparent, verifiable, and
accessible to others. In an era where the scientific community faces
concerns about the replicability of research findings, adopting
reproducible research practices becomes imperative.

This document uses Allaire (2023).

\section{Key Benefits of Reproducible
Research}\label{key-benefits-of-reproducible-research}

\begin{itemize}
\tightlist
\item
  \emph{Increased Credibility}: Transparent and reproducible research
  enhances the credibility of scientific findings, as others can
  independently verify and validate the results.
\item
  \emph{Collaboration and Iteration}: Open sharing of code and data
  facilitates collaboration among researchers and allows for the
  iterative improvement of studies over time.
\item
  \emph{Efficient Communication}: Reproducible research documents serve
  as comprehensive and self-contained artifacts, aiding in effective
  communication of methodologies and results.
\end{itemize}

\section{Method}\label{method}

\subsection{Getting started}\label{getting-started}

\begin{enumerate}
\def\labelenumi{\arabic{enumi}.}
\tightlist
\item
  Install \href{https://quarto.org/docs/get-started/}{Quarto}.
\item
  You may need to install Public Sans font if you do not have it. You
  can find it
  \href{https://github.com/anuopensci/quarto-anu-report/tree/main/_extensions/anu-report/assets/webfonts}{here}.
\item
  The easiest way to edit and render the document is using the
  \href{https://posit.co/download/rstudio-desktop/}{RStudio IDE}. If you
  prefer not to install R, then you can use VS Code instead.
\item
  From the Terminal run the following command
\end{enumerate}

\begin{Shaded}
\begin{Highlighting}[]
\ExtensionTok{quarto}\NormalTok{ use template anuopensci/quarto{-}anu{-}report}
\end{Highlighting}
\end{Shaded}

\section{Conclusion}\label{conclusion}

This report underscores the transformative impact of reproducible
research on the scientific landscape, promoting a commitment to openness
and accountability for the benefit of the entire research community.

\section*{References}\label{references}
\addcontentsline{toc}{section}{References}

\phantomsection\label{refs}
\begin{CSLReferences}{1}{0}
\bibitem[\citeproctext]{ref-quarto}
Allaire, JJ. 2023. \emph{Quarto: R Interface to 'Quarto' Markdown
Publishing System}. \url{https://CRAN.R-project.org/package=quarto}.

\end{CSLReferences}


\backmatter


\end{document}
